\documentclass[11pt]{article}
\usepackage[T1]{fontenc}
\usepackage{lmodern}
\usepackage{amsmath,amssymb,amsthm}
\usepackage{graphicx}
\usepackage{hyperref}
\usepackage[margin=1in]{geometry}

\newtheorem{theorem}{Theorem}
\newtheorem{lemma}{Lemma}
\newtheorem{remark}{Remark}

\newcommand{\R}{\mathbb{R}}
\newcommand{\N}{\mathbb{N}}
\newcommand{\clco}{\overline{\operatorname{co}}}

\title{A weakly closed counterexample for remotality of convex hulls}
\author{arXiv automatic math run \texttt{fa\_banach\_001}, agent \texttt{agent\_super\_01}}
\date{July 1, 2026}

\begin{document}
\maketitle

\section*{Status}

This packet records a candidate counterexample, likely valid, to a question in
Martin--Rao \cite{MR}.  The source paper asks whether, for a weakly closed and
bounded set \(E\), remotality of \(K=\clco(E)\) from a point forces remotality
of \(E\) from that point.  It also asks whether the answer is positive for RNP
sets.  The construction below gives a negative answer to both versions.

\section*{Source question}

The question appears in Remark 8 of \cite{MR}, immediately after an example
showing that norm closedness is not enough:
\begin{quote}
Let \(X\) be a Banach space and let \(E \subset X\) be a weakly closed and
bounded set. We do not know if remotality of \(K=\clco(E)\) from a point always
implies that of \(E\). [...] We may also ask whether the above question has a
positive answer for RNP sets.
\end{quote}

The source paper's main theorem solves the headline problem of existence of
closed bounded convex non-remotal sets in every infinite-dimensional Banach
space; the remark above is a separate residual question about passing
remotality from a closed convex hull back to its generating set.

\begin{figure}[h]
\centering
\includegraphics[width=.88\textwidth]{figures/open_problem_crop.png}
\caption{The source page containing Remark 8 of arXiv:0909.1992.}
\end{figure}

\section*{Definitions}

A bounded set \(A\subset X\) is \emph{remotal from \(x\in X\)} if there is
\(a_0\in A\) such that
\[
  \|x-a_0\|=\sup\{\|x-a\|:a\in A\}.
\]
It is enough to consider \(x=0\), since the general case is obtained by
translation.

We use \(K=\clco(E)\) for the norm-closed convex hull.  A closed bounded convex
set is an RNP set if it has the Radon--Nikodym property in the standard convex
set sense; in particular every closed bounded convex subset of a Banach space
with the Radon--Nikodym property is an RNP set.

\section*{Proof intuition}

Martin--Rao's norm-closed example fails to be weakly closed because its Hilbert
basis coordinate converges weakly to zero, and the missing weak limit is
exactly the farthest point created in the convex hull.  The fix is to attach a
second, smaller ``tag'' coordinate which prevents weak convergence of the
individual generators, but which can still be cancelled by taking signed
midpoints.

The summing basis \(s_n=(1,\ldots,1,0,\ldots)\) in \(c_0\) has precisely this
behavior: it has no weak limit in \(c_0\), since coordinate functionals would
force any limit to be the constant sequence \(1\), but a signed pair
\(\pm s_n\) cancels exactly.  Pairing this tag with the usual Hilbert basis
coordinate gives a weakly closed set \(E\); averaging each signed pair gives
\((a_n,0,0)\), and these averages converge in norm to a unit vector in
\(\clco(E)\).

\section*{The counterexample}

\begin{theorem}
There is a Banach space \(X\), a weakly closed bounded set \(E\subset X\), and
\(K=\clco(E)\), such that \(K\) is remotal from \(0\) but \(E\) is not remotal
from \(0\).  Moreover, \(K\) may be chosen to be an RNP set.
\end{theorem}

\begin{proof}
Let
\[
  X=\R\oplus_\infty \ell_2\oplus_\infty c_0
\]
with norm
\[
  \|(t,h,z)\|=\max\{|t|,\|h\|_2,\|z\|_\infty\}.
\]
Let \((e_n)\) be the canonical orthonormal basis of \(\ell_2\), and let
\[
  s_n=(\underbrace{1,\ldots,1}_{n\text{ times}},0,0,\ldots)\in c_0
\]
be the summing basis of \(c_0\).  Put \(a_n=n/(n+1)\) and fix \(c=1/2\).  Define
\[
  x_n^+=(a_n,a_ne_n,cs_n),\qquad
  x_n^-=(a_n,-a_ne_n,-cs_n)
\]
and
\[
  E=\{x_n^+,x_n^-:n\in\N\}.
\]

The set \(E\) is bounded and every element of \(E\) has norm
\[
  \|x_n^\pm\|=\max\{a_n,a_n,c\}=a_n<1.
\]
Hence \(E\) has no vector of maximal norm, while
\[
  \sup_{x\in E}\|x\|=\sup_n a_n=1.
\]
Thus \(E\) is not remotal from \(0\).

We next prove that \(E\) is weakly closed.  Suppose a net
\((x_{n_i}^{\sigma_i})_i\), with \(\sigma_i\in\{+,-\}\), converges weakly in
\(X\) to some point \(x\).  If the indices \(n_i\) are eventually contained in
a finite subset of \(\N\), then a subnet is eventually constant and the weak
limit belongs to the finite set \(E\).  So it remains to rule out the case in
which \(n_i\) tends to infinity along a subnet.

Passing to this subnet, the \(c_0\)-coordinates
\[
  \sigma_i c s_{n_i}
\]
would converge weakly in \(c_0\).  Evaluating at the first coordinate
functional \(\delta_1\in \ell_1=c_0^*\), the scalars
\(\delta_1(\sigma_i c s_{n_i})=\sigma_i c\) must converge.  Since they take
only the two values \(\pm c\), they are eventually constant; write the eventual
sign as \(\sigma\).  For every fixed \(k\), once \(n_i\geq k\) we have
\[
  \delta_k(\sigma_i c s_{n_i})=\sigma c.
\]
Therefore a weak limit \(z\in c_0\) would satisfy \(z_k=\sigma c\) for every
\(k\), impossible because \(c>0\) and the constant nonzero sequence is not in
\(c_0\).  Thus no such subnet exists, and \(E\) is weakly closed.

Now let \(K=\clco(E)\).  For each \(n\),
\[
  \frac{x_n^++x_n^-}{2}=(a_n,0,0)\in \operatorname{co}(E).
\]
Since \(a_n\to 1\), the point
\[
  u=(1,0,0)
\]
belongs to \(K\).  Also \(E\subset B_X\), so \(K\subset B_X\).  Hence
\[
  \sup_{y\in K}\|y\|=1=\|u\|,
\]
and \(K\) is remotal from \(0\).

Finally, \(X^*=\R\oplus_1\ell_2\oplus_1\ell_1\) is separable.  Therefore \(X\)
has the Radon--Nikodym property.  Since \(K\) is a closed bounded convex subset
of \(X\), it is an RNP set.  This gives the asserted negative answer even for
RNP sets.
\end{proof}

\section*{Verification and novelty checks}

The proof is deterministic and uses only elementary properties of
\(\ell_2\), \(c_0\), and the Radon--Nikodym property.  The key point is
weak closedness of \(E\), supplied by the summing-basis tag in \(c_0\).

Before promotion, the following bounded novelty checks were performed.
The run indexes
\[
\texttt{registry\_index.tsv},\quad \texttt{solutions/index.tsv},\quad
\texttt{attempts/index.tsv},\quad \texttt{proof\_gaps/index.tsv}
\]
were searched for arXiv:0909.1992, remotality, remotal sets, farthest points,
closed convex hulls, weakly closed bounded sets, and RNP-set variants.  No
existing answer to this source question was found.  The local parsed arXiv
corpus was searched for remotality, farthest point, weakly closed, closed
convex hull, and Martin--Rao/Sababheh--Khalil phrases; the visible later
material referenced the main Martin--Rao theorem but did not answer Remark 8.
External web searches used exact phrases such as ``We do not know if
remotality'' with ``cconv(E)'', ``remotality'' with ``weakly closed'' and
``closed convex hull'', and ``remotal'' with ``RNP'' and ``closed convex
hull''.  These searches
found no later solution to the specific weakly closed/RNP convex-hull question.

\section*{Human-review recommendation}

Review as a likely valid counterexample.  The two points most worth checking
are:
\begin{enumerate}
  \item the net argument proving that \(E\) is weakly closed;
  \item the interpretation of the source phrase ``for RNP sets''.  The packet
  proves the stronger-looking formulation where \(K=\clco(E)\) is a closed
  bounded convex subset of a Banach space with the Radon--Nikodym property.
\end{enumerate}

\begin{thebibliography}{9}

\bibitem{MR}
M. Martin and T. S. S. R. K. Rao,
\emph{On remotality for convex sets in Banach spaces},
Journal of Approximation Theory 162 (2010), 392--396;
arXiv:0909.1992.

\end{thebibliography}

\end{document}
